\section{2025-02-13}

\startplacetable[location={here},title={Bin vs Normal}]
\setupTABLE[r][1][background=color,backgroundcolor=gray]
\bTABLE
\startCSV
Binominal, Normal
Diskret, Stetig
[nc=2]Wahrscheinlichkeitsverteilungen
, Stetigkeitskorrektur
\stopCSV
\eTABLE
\stopplacetable

Varianz: Mittelwert der quadratischen Abweichung? $n * p * q$


Erwartungswert: Mittelwert an den man sich bei genug versuchen annähert

\startframedtipp
starte bei dem offensichtlichsten
\stopframedtipp


\startitemize[packed]
\item Laplace: Alle Ergeignisse sind gleich wahrscheinlich
\stopitemize


\startframedtipp
\startitemize[packed]
\item Bedingte Wahrscheilichkeiten
\item Zwei Fragen zu einem sachverhalt, die gleich erscheinen unterscheiden sich wahrscheinlich in einer bedingten Wahrscheinlichkeit
\stopitemize
\stopframedtipp


\section{2025-02-17}

\startframedtipp
\startitemize[packed]
\item In Teil A (Hilfsmittel) auch, wennn man nicht weiter kommt, mal in Teil B gucken, vielleicht sind da anreize?
\item Beim Vergleichen: {\bf Auch} die Einfachen / Offensichtlichen Aspekte hinschreiben
\item Es kann auch mal eine \quotation{einfache} Aufgabe drann kommen
\stopitemize
\stopframedtipp


\section{2025-03-25}

\startitemize[packed]
\item Aufgaben mathematisieree
\item Bei \quotation{Bestimmen Sie den höchsten Stromverbrauch, wenn der
Stromverbrauch mit der Funktion $f(x)$ beschrieben werden kann}, reicht als
Antwort \quotation{Da das Maximum der Funktion bei $x_1$ liegt, ist dort auch
der höchste Stromverbrauch mit $y_1$}. Wenn das Maximum bestimmt werden soll,
{\bf reicht das nicht}. Man kann nicht etwas (das Maximum) mit sich selber (dem
Maximum) bestimmen.
\stopitemize
